\section{Основная теорема бесконечной теории Галуа}

\begin{proposition}[7.12]
Пусть $\Om/F$~--- расширение Галуа с группой $G$.
\begin{enumerate}[label=\textup{(\alph*)}]
\item Если $F\subset M\subset\Om$, то $\Om/M$~--- расширение Галуа, подгруппа $\Gal(\Om/M)$ замкнута в $G$ и
$\Om^{\Gal(\Om/M)}=M$.
\item Для любой подгруппы $H\subset G$ подгруппа $\Gal(\Om/\Om^H)$ равна замыканию $\bar H$ подгруппы $H$.
\end{enumerate}
\end{proposition}

\begin{proof}
(a) Первое утверждение~--- \rf{7.3}. Далее, $\Gal(\Om/M)=\bigcap_{S}G(S)$ по всем конечным $S\subset M$; каждая $G(S)$
открыта, а значит, замкнута (\rf{7.1a}), и пересечение замкнутых множеств замкнуто. Последнее утверждение~--- \rf{7.9} для
расширения $\Om/M$.

(b) Подгруппа $\Gal(\Om/\Om^H)$ содержит $H$ и замкнута по (a), значит, содержит $\bar H$. Обратно, пусть
$\sigma\in G\setminus\bar H$; покажем, что $\sigma$ двигает некоторый элемент $\Om^H$. Так как $\sigma\notin\bar H$, а
множества $\sigma\Gal(\Om/E)$ ($E$~--- конечное подрасширение Галуа) образуют базу окрестностей $\sigma$, найдётся $E$
с $\sigma\Gal(\Om/E)\cap H=\varnothing$. Пусть $\pi_E\colon G\to\Gal(E/F)$~--- ограничение. Тогда
$\sigma|_E\notin\pi_E(H)$: иначе $\sigma|_E=h|_E$ для некоторого $h\in H$, т.\,е.\ $h\in\sigma\Gal(\Om/E)$. По основной
теореме для конечного расширения $E/F$ (или по теореме Артина \rf{3.11}) $\Gal(E/E^{\pi_E(H)})=\pi_E(H)$, поэтому $\sigma|_E$
двигает некоторый $x\in E^{\pi_E(H)}$. Но $E^{\pi_E(H)}\subset\Om^H$: для $h\in H$ имеем $h(x)=h|_E(x)=x$. Итак,
$\sigma\notin\Gal(\Om/\Om^H)$.
\end{proof}

\begin{center}
\begin{tikzpicture}[scale=0.85]
\draw[cold,fill=cold!8] plot[smooth cycle] coordinates {(-2,0) (-1.2,1.1) (0.6,1.0) (1.2,0) (0.3,-0.9) (-1.4,-0.9)};
\node[cold] at (-1.75,1.15) {$\bar H$};
\foreach \p in {(-1.3,0.2),(-0.8,-0.35),(-0.3,0.5),(0.2,-0.2),(0.6,0.4),(-1.0,0.7)} \fill[cold] \p circle (1.3pt);
\node[cold,font=\footnotesize] at (-0.35,-0.65) {точки $H$};
\fill[cnew] (3.0,0.2) circle (1.6pt) node[below right]{$\sigma$};
\draw[cnew,dashed] (3.0,0.2) circle (0.75);
\node[cnew,font=\footnotesize,align=left,anchor=west] at (4.0,0.2) {$\sigma\Gal(\Om/E)$~--- окрестность $\sigma$,\\ не пересекающая $H$};
\end{tikzpicture}
\end{center}
\figcap{Если $\sigma\notin\bar H$, то $\sigma$ отличим от всех элементов $H$ уже на конечном уровне $E$:
$\sigma|_E\notin\pi_E(H)$.}

\begin{theorem}[7.13: основная теорема]
Пусть $\Om/F$~--- расширение Галуа с группой $G$. Отображения
\[
H\mapsto\Om^H,\qquad M\mapsto\Gal(\Om/M)
\]
~--- взаимно обратные биекции между множеством замкнутых подгрупп $G$ и множеством промежуточных полей $F\subset
M\subset\Om$:
\[
\{\text{замкнутые подгруппы }G\}\ \overset{1:1}{\longleftrightarrow}\ \{\text{промежуточные поля }F\subset M\subset\Om\}.
\]
При этом:
\begin{enumerate}[label=\textup{(\alph*)}]
\item $H_1\supset H_2\iff\Om^{H_1}\subset\Om^{H_2}$ (соответствие обращает включения);
\item замкнутая подгруппа $H$ открыта $\iff$ $[\Om^H:F]<\infty$; в этом случае $(G:H)=[\Om^H:F]$;
\item $\sigma H\sigma^{-1}\leftrightarrow\sigma M$, т.\,е.\ $\Om^{\sigma H\sigma^{-1}}=\sigma(\Om^H)$ и
$\Gal(\Om/\sigma M)=\sigma\Gal(\Om/M)\sigma^{-1}$;
\item замкнутая подгруппа $H$ нормальна $\iff$ $\Om^H/F$~--- расширение Галуа; в этом случае ограничение
$\sigma\mapsto\sigma|_{\Om^H}$ индуцирует изоморфизм \emph{топологических} групп $G/H\cong\Gal(\Om^H/F)$ (на $G/H$~---
фактор-топология).
\end{enumerate}
\end{theorem}

\begin{proof}
\emph{Взаимная обратность.} Если $H$ замкнута, то $\Gal(\Om/\Om^H)=\bar H=H$ по \rf{7.12}(b). Если $M$~--- промежуточное
поле, то $\Gal(\Om/M)$ замкнута и $\Om^{\Gal(\Om/M)}=M$ по \rf{7.12}(a).

(a) Импликации $H_1\supset H_2\Rightarrow\Om^{H_1}\subset\Om^{H_2}\Rightarrow\Gal(\Om/\Om^{H_1})\supset\Gal(\Om/\Om^{H_2})$
очевидны, а последние группы равны $H_1$ и $H_2$.

(b) Замкнутая подгруппа конечного индекса открыта, а открытая подгруппа компактной группы имеет конечный индекс
(\rf{7.1a}). Итак, для замкнутой подгруппы $H$ компактной группы
\[
\boxed{\ H\ \text{открыта}\iff(G:H)<\infty.\ }
\] Пусть $H$ открыта. Тогда $H$ содержит базисную окрестность $\Gal(\Om/K)$, $K$~--- конечное подрасширение Галуа,
и по (a) $\Om^H\subset K$, так что $[\Om^H:F]<\infty$. Обратно, пусть $M=\Om^H$ конечно над $F$, $M=F(x_1,\dots,x_n)$. Тогда
$H=\Gal(\Om/M)=G(\{x_1,\dots,x_n\})$~--- открытая подгруппа по определению топологии Крулля.

Осталось доказать $(G:H)=[M:F]$ для $M=\Om^H$ конечного над $F$. Здесь $G/H$~--- множество \emph{левых смежных классов}
(подгруппа $H$ не обязана быть нормальной). Отображение
\[
G/H\to\Hom_F(M,\Om),\qquad\sigma H\mapsto\sigma|_M,
\]
корректно (элементы $H$ тождественны на $M$), инъективно (если $\sigma|_M=\tau|_M$, то $\tau^{-1}\sigma\in\Gal(\Om/M)=H$)
и сюръективно (каждый $F$-гомоморфизм $M\to\Om$ продолжается до элемента $G$ по \rf{7.4}). Значит, $(G:H)=|\Hom_F(M,\Om)|$.
Расширение $M/F$ конечно и сепарабельно, поэтому $M=F(\alpha)$ (теорема о примитивном элементе \rf{5.1}); $F$-гомоморфизмы
$F(\alpha)\to\Om$ взаимно однозначно соответствуют корням минимального многочлена $\mu_\alpha$ в $\Om$, а их ровно
$\deg\mu_\alpha=[M:F]$, так как $\mu_\alpha$ сепарабелен и раскладывается в $\Om$ (нормальность).
\[
\begin{tikzcd}
M=\Om^H\arrow[r,hook]\arrow[dr,"\varphi"'] & \Om\arrow[d,dashed,cnew,"\sigma\ (7.4)"]\\
 & \Om
\end{tikzcd}
\qquad\qquad
\sigma H\ \longleftrightarrow\ \varphi=\sigma|_M
\]
\emph{Другой взгляд на конечность.} Биекция $G/H\to\Hom_F(\Om^H,\Om)$ построена без предположения
$[\Om^H:F]<\infty$. Поэтому для открытой $H$ множество $\Hom_F(\Om^H,\Om)$ конечно. Если бы $[\Om^H:F]=\infty$, то в
$\Om^H$ нашлись бы конечные подрасширения $K'$ сколь угодно большой степени $d$ (иначе подрасширение $K$ наибольшей
степени совпадало бы с $\Om^H$: для $x\in\Om^H$ степень $K(x)$ не больше степени $K$, т.\,е.\ $K(x)=K$). У каждого такого
$K'$ ровно $d$ различных $F$-вложений в $\Om$, и по \rf{7.4} каждое продолжается до вложения $\Om^H\to\Om$; получилось бы
$|\Hom_F(\Om^H,\Om)|\ge d$ для всех $d$.

(c) Для $\tau\in G$ и $\alpha\in\Om$: $\tau\alpha=\alpha\iff\sigma\tau\sigma^{-1}(\sigma\alpha)=\sigma\alpha$. Поэтому
$\Gal(\Om/\sigma M)=\sigma\Gal(\Om/M)\sigma^{-1}$, и по взаимной обратности $\Om^{\sigma H\sigma^{-1}}=\sigma(\Om^H)$.

(d) Пусть $M=\Om^H$. По (c) и взаимной обратности $H$ нормальна $\iff$ $\sigma M=M$ для всех $\sigma\in G$. Если
$\sigma M=M$ для всех $\sigma$, то $M/F$~--- Галуа по \rf{7.5}. Обратно, если $M/F$~--- Галуа, то для $\alpha\in M$ элемент
$\sigma(\alpha)$~--- корень минимального многочлена $\alpha$, все корни которого лежат в $M$; значит, $\sigma M\subset M$
(и, применяя к $\sigma^{-1}$, $\sigma M=M$).

Пусть теперь $H$ нормальна. Ограничение $r\colon G\to\Gal(M/F)$, $\sigma\mapsto\sigma|_M$, определено ($\sigma M=M$),
является гомоморфизмом, сюръективно (по \rf{7.4} каждый элемент $\Gal(M/F)$ продолжается до элемента $G$) и имеет ядро
$\Gal(\Om/M)=H$. Значит, $r$ индуцирует изоморфизм групп $\bar r\colon G/H\to\Gal(M/F)$.
\[
\begin{tikzcd}[column sep=large]
G\arrow[r,two heads,"\pi"]\arrow[rr,bend left=22,"{r\colon\ \sigma\,\mapsto\,\sigma|_M}"] & G/H\arrow[r,cnew,"\bar r","\sim"'] & \Gal(M/F)
\end{tikzcd}
\] Он непрерывен: базу
окрестностей единицы в $\Gal(M/F)$ образуют подгруппы $\Gal(M/K)$, $K\subset M$~--- конечное подрасширение Галуа, а
$r^{-1}(\Gal(M/K))=\Gal(\Om/K)$ открыта в $G$; значит, $r$ непрерывен (\rf{7.7a}), а $\bar r$ непрерывен по определению
фактор-топологии. (явно: $\bar r^{-1}(\Gal(M/K))=\pi(\Gal(\Om/K))$, и это множество открыто в
$G/H$, так как его прообраз $\pi^{-1}\pi(\Gal(\Om/K))=\Gal(\Om/K)\,H$~--- объединение открытых сдвигов $\Gal(\Om/K)h$). Группа $G/H$ компактна (непрерывный образ компакта $G$), $\Gal(M/F)$ хаусдорфова (\rf{7.8}), а непрерывная
биекция компакта на хаусдорфово пространство~--- гомеоморфизм.
\end{proof}

\begin{remark}[7.14]
Как и в конечном случае (\rf{3.18}):
\begin{enumerate}[label=(\alph*),leftmargin=*]
\item Пусть $(M_i)_{i\in I}$~--- семейство (возможно бесконечное) промежуточных полей и $H_i\leftrightarrow M_i$. Пусть
$M$~--- композит $M_i$ (наименьшее поле, содержащее все $M_i$). Так как $\bigcap_iH_i$~--- наибольшая (замкнутая)
подгруппа, содержащаяся во всех $H_i$, то $\Gal(\Om/M)=\bigcap_iH_i$.
\item Пусть $M\leftrightarrow H$. Наибольшая (замкнутая) нормальная подгруппа, содержащаяся в $H$,~--- это
$N=\bigcap_{\sigma\in G}\sigma H\sigma^{-1}$, поэтому $\Om^N$~--- композит полей $\sigma M$~--- является наименьшим
нормальным расширением $F$, содержащим $M$ (нормальным замыканием $M$ в $\Om$).
\end{enumerate}
\end{remark}

\begin{aside}[7.14a: зачем нужна топология]
Отображение $H\mapsto\Om^H$ определено для \emph{всех} подгрупп, но различает лишь их замыкания: $\Om^H=\Om^{\bar H}$
(\rf{7.12}(b)). Если ограничиться абстрактными подгруппами, соответствие Галуа перестаёт быть биекцией, и это не
<<патология на краю теории>>, а типичная ситуация.
\begin{itemize}[leftmargin=*]
\item В $\Gal(\Fpb/\F_p)$ подгруппа $\langle\Frob\rangle\cong\Z$, порождённая Фробениусом, собственная, но её
неподвижное поле~--- $\F_p$, как и у всей группы (\rf{7.17}).
\item В $\Gal(\Qal/\Q)$ есть подгруппы индекса~2, неподвижное поле которых~--- $\Q$ (\rf{7.29}, \rf{7.30}). Даже <<конечные расширения
$\leftrightarrow$ подгруппы конечного индекса>> неверно: конечным расширениям соответствуют \emph{открытые} подгруппы.
\end{itemize}
Топология отбирает среди подгрупп ровно те, что <<видны>> полям. Так же среди гомоморфизмов группы Галуа в топологические группы
(например, в $\mathrm{GL}_n(\CC)$ или в $\mathrm{GL}_n(\Q_\ell)$) естественны лишь непрерывные: их значения на
элементе с любой точностью определяются его действием на достаточно большом конечном подрасширении. Поэтому в теории чисел представления групп Галуа по определению предполагаются непрерывными.
\end{aside}

\begin{proposition}[7.15]
Пусть $E$ и $L$~--- расширения $F$, содержащиеся в общем поле. Если $E/F$~--- расширение Галуа, то $EL/L$ и
$E/E\cap L$~--- расширения Галуа, и ограничение
\[
\sigma\mapsto\sigma|_E\colon\Gal(EL/L)\to\Gal(E/E\cap L)
\]
~--- изоморфизм топологических групп.
\end{proposition}

\begin{center}
\begin{tikzpicture}[scale=0.9]
\node (EL) at (0,2.6) {$EL$};
\node (E) at (-1.8,1.3) {$E$};
\node (L) at (1.8,1.3) {$L$};
\node (I) at (0,0) {$E\cap L$};
\node (F) at (0,-1.3) {$F$};
\draw (E) -- (EL); \draw (I) -- (L); \draw (I) -- (F);
\draw[cnew,thick] (I) -- node[above right,font=\footnotesize,pos=0.45]{Галуа} (E);
\draw[cnew,thick] (L) -- node[above right,font=\footnotesize]{Галуа} (EL);
\draw (F) to[out=180,in=-100] node[left,font=\footnotesize,pos=0.4]{Галуа} (E);
\end{tikzpicture}
\end{center}
\figcap{Параллельные красные стороны <<ромба>> имеют изоморфные группы Галуа: $\Gal(EL/L)\cong\Gal(E/E\cap L)$.}

\begin{proof}
$E/E\cap L$~--- Галуа по \rf{7.3}. Расширение $EL/L$ порождено элементами $E$; минимальный многочлен над $L$ элемента $x\in E$
делит его минимальный многочлен над $F$, который сепарабелен и раскладывается в $E\subset EL$. Значит, $EL$~---
объединение полей разложения сепарабельных многочленов над $L$, т.\,е.\ $EL/L$~--- Галуа. Пусть
$G_1=\Gal(EL/L)$, $G_2=\Gal(E/E\cap L)$, $\varphi(\sigma)=\sigma|_E$.
\begin{itemize}[leftmargin=*]
\item \emph{Корректность:} $\sigma(E)=E$, так как $E/F$ нормально, и $\sigma$ тождествен на $E\cap L\subset L$.
\item \emph{Инъективность:} если $\sigma$ тождествен на $E$ и на $L$, то и на $EL$.
\item \emph{Непрерывность:} для конечного $S\subset E$ прообраз $G_2(S)$ равен $G_1(S)$ (\rf{7.7a}).
\item \emph{Сюръективность.} Образ $\varphi(G_1)$ компактен, значит, замкнут в $G_2$ (\rf{7.8}). Его неподвижное поле:
если $x\in E$ и $\sigma(x)=x$ для всех $\sigma\in G_1$, то $x\in(EL)^{G_1}=L$ (\rf{7.9}), т.\,е.\ $x\in E\cap L$. Итак,
$E^{\varphi(G_1)}=E\cap L$, и по \rf{7.13} $\varphi(G_1)=\Gal(E/E\cap L)$.
\item \emph{Гомеоморфизм:} непрерывная биекция компакта на хаусдорфово пространство.\qedhere
\end{itemize}
\end{proof}

\begin{lemma}[7.15a]
Пусть $E/F$~--- нормальное алгебраическое расширение, $\Om$~--- поле, содержащее $F$, и $\varphi,\psi\colon E\to\Om$~---
$F$-гомоморфизмы. Тогда $\varphi(E)=\psi(E)$.
\end{lemma}

\begin{proof}
Покажем, что $\varphi(E)=R$, где $R\subset\Om$~--- множество корней в $\Om$ минимальных многочленов элементов $E$ (оно не
зависит от $\varphi$). Если $\alpha\in E$, то $\mu_\alpha(\varphi\alpha)=\varphi(\mu_\alpha(\alpha))=0$, так что
$\varphi(E)\subset R$. Обратно, пусть $\beta\in\Om$~--- корень $\mu_\alpha$, $\alpha\in E$. По нормальности
$\mu_\alpha=(X-\alpha_1)\cdots(X-\alpha_d)$ с $\alpha_i\in E$. Применяя $\varphi$ к коэффициентам (они в $F$), получаем
$\mu_\alpha=(X-\varphi\alpha_1)\cdots(X-\varphi\alpha_d)$ в $\Om[X]$; корни $\mu_\alpha$ в $\Om$~--- ровно
$\varphi\alpha_1,\dots,\varphi\alpha_d$, так что $\beta=\varphi(\alpha_i)\in\varphi(E)$.
\end{proof}

\begin{corollary}[7.16]
Пусть $\Om$~--- алгебраически замкнутое поле, содержащее $F$, а $E$ и $L$~--- как в \rf{7.15}. Если $\rho\colon E\to\Om$ и
$\sigma\colon L\to\Om$~--- $F$-гомоморфизмы и $\rho|_{E\cap L}=\sigma|_{E\cap L}$, то существует $F$-гомоморфизм
$\tau\colon EL\to\Om$ с $\tau|_E=\rho$ и $\tau|_L=\sigma$.
\end{corollary}

\[
\begin{tikzcd}[row sep=small,column sep=large]
 & E\arrow[dr,hook]\arrow[drr,bend left=25,"\rho"] & & \\
E\cap L\arrow[ur,hook]\arrow[dr,hook] & & EL\arrow[r,dashed,cnew,"\tau"] & \Om\\
 & L\arrow[ur,hook]\arrow[urr,bend right=25,"\sigma"'] & &
\end{tikzcd}
\]
\figcap{Склейка: вложения $\rho$ и $\sigma$, согласованные на $E\cap L$, склеиваются во вложение композита (красная
стрелка).}

\begin{proof}
\emph{1)} Расширение $EL/L$ алгебраично, поэтому $\sigma$ продолжается до гомоморфизма $s\colon EL\to\Om$ (теорема \rf{6.6},
применённая к полю $L$ и алгебраическому замыканию поля $\sigma(L)$ в $\Om$):
\[
\begin{tikzcd}
L\arrow[r,hook]\arrow[dr,"\sigma"'] & EL\arrow[d,dashed,cnew,"s"]\\
 & \Om
\end{tikzcd}
\]

\emph{2)} Отображения $s|_E$ и $\rho$~--- $F$-вложения $E\to\Om$, совпадающие на $E\cap L$ (оба равны там $\sigma$). По
лемме \rf{7.15a} $s(E)=\rho(E)$, поэтому $\varepsilon=\rho^{-1}\circ s|_E\colon E\to E$~--- автоморфизм (инъективен как
гомоморфизм полей, сюръективен, так как $s(E)=\rho(E)$), тождественный на $E\cap L$. Итак, $\varepsilon\in\Gal(E/E\cap L)$
и $s|_E=\rho\circ\varepsilon$.

\emph{3)} По \rf{7.15} существует (единственный) $e\in\Gal(EL/L)$ с $e|_E=\varepsilon$. Положим $\tau=s\circ e^{-1}$. Так
как $e$ тождествен на $L$, то $\tau|_L=s|_L=\sigma$. Так как $e(E)=E$ и $e^{-1}|_E=\varepsilon^{-1}$, то
$\tau|_E=s|_E\circ\varepsilon^{-1}=\rho$:
\[
\begin{tikzcd}[column sep=large,row sep=large]
E\arrow[r,hook]\arrow[d,cnew,"\varepsilon"'] & EL\arrow[r,"s"]\arrow[d,"e"'] & \Om\\
E\arrow[r,hook] & EL\arrow[ur,cnew,"\tau=s\circ e^{-1}"'] &
\end{tikzcd}
\]
\end{proof}

\begin{remark}[7.16a]
Нормальность $E$ здесь существенна. Если $E=\Q(\sqrt[3]2)$ и $L=\Q(\omega\sqrt[3]2)$ ($\omega=e^{2\pi i/3}$), то
$E\cap L=\Q$, но вложения $\sqrt[3]2\mapsto\sqrt[3]2$ и $\omega\sqrt[3]2\mapsto\sqrt[3]2$ не склеиваются: иначе
$\omega=\omega\sqrt[3]2/\sqrt[3]2\mapsto1$.
\end{remark}

Следствие \rf{7.16}~--- это \emph{принцип склейки}: вложение композита задаётся парой вложений, согласованных на пересечении.
Его главное применение~--- описание групп Галуа композитов, в том числе бесконечных, где подсчёт порядков невозможен.

\begin{corollary}[7.16b]
Пусть $E$ и $L$~--- расширения Галуа поля $F$ (возможно бесконечные) в общем поле. Тогда $EL/F$~--- расширение Галуа и
\[
\Gal(EL/F)\to\{(a,b)\in\Gal(E/F)\times\Gal(L/F)\mid a|_{E\cap L}=b|_{E\cap L}\},\qquad\sigma\mapsto(\sigma|_E,\sigma|_L),
\]
~--- изоморфизм топологических групп (справа~--- \emph{расслоенное произведение} $\Gal(E/F)\times_{\Gal(E\cap L/F)}\Gal(L/F)$
с топологией подпространства). В частности, если $E\cap L=F$, то $\Gal(EL/F)\cong\Gal(E/F)\times\Gal(L/F)$.
\end{corollary}

\begin{proof}
$EL$~--- объединение полей разложения сепарабельных многочленов над $F$, так что $EL/F$~--- Галуа. Расширение
$E\cap L/F$ тоже Галуа: минимальный многочлен над $F$ элемента $E\cap L$ сепарабелен и раскладывается и в $E$, и в $L$,
т.\,е.\ в $E\cap L$. Отображение
определено ($E$, $L$ нормальны), инъективно ($EL$ порождено $E$ и $L$) и непрерывно (обе компоненты~--- ограничения,
непрерывные по \rf{7.7a}). Сюръективность: пусть $a|_{E\cap L}=b|_{E\cap L}$. Возьмём алгебраическое замыкание
$\Om\supset EL$ и применим \rf{7.16} к $\rho=a$, $\sigma=b$: получим $\tau\colon EL\to\Om$ с $\tau|_E=a$, $\tau|_L=b$; так как
$\tau(EL)$ порождено $a(E)=E$ и $b(L)=L$, то $\tau\in\Gal(EL/F)$. Расслоенное произведение замкнуто в хаусдорфовом
произведении (это множество, где совпадают два непрерывных отображения в хаусдорфову группу $\Gal(E\cap L/F)$);
гомеоморфность~--- по компактности, как выше.
\end{proof}

Для конечных $E$, $L$ это \rf{3.21}; например, для $E=\Q(\sqrt[3]2,\omega)$ и $L=\Q(\sqrt{-3})$ имеем $L\subset E$, и
расслоенное произведение равно $\Gal(E/\Q)\cong S_3$, а не $S_3\times\Z/2$. Бесконечные примеры~--- в \rf{7.36}.

\begin{example}[7.17: алгебраическое замыкание конечного поля]
Пусть $\Om=\Fpb$~--- алгебраическое замыкание $\F_p$ и $G=\Gal(\Om/\F_p)$. Поле $\F_p$ совершенно, так что $\Fpb=\F_p^{\rm sep}$ и $\Om/\F_p$~--- расширение Галуа. Напомним (\rf{4.20}, \rf{4.23}): $\Om$
содержит ровно одно подполе $\F_{p^n}$ из $p^n$ элементов, $\F_{p^m}\subset\F_{p^n}\iff m\mid n$, и
$\Gal(\F_{p^n}/\F_p)$~--- циклическая группа порядка $n$, порождённая $\sigma_n\colon a\mapsto a^p$. Конечные
подрасширения $\Om$~--- это ровно поля $\F_{p^n}$, и все они Галуа.

\emph{Фробениус порождает $G$ как топологическую группу.} Пусть $\sigma\colon a\mapsto a^p$~--- элемент Фробениуса в $G$
и $H=\langle\sigma\rangle$. Тогда $\Om^H=\{x\in\Om\mid x^p=x\}=\F_p$ (у $X^p-X$ не больше $p$ корней), и по \rf{7.12}(b)
$\bar H=\Gal(\Om/\Om^H)=G$.

\emph{Топология на $\Z$.} Подгруппы $n\Z$ ($n\ge1$) удовлетворяют условиям \rf{7.2} (a)--(d): $m\Z\cap n\Z=\mathrm{НОК}(m,n)\Z$,
а $\Z$ абелева. Получаем топологию на $\Z$, в которой $\{n\Z\}$~--- база окрестностей нуля (\emph{проконечная топология}).
Открытые множества~--- объединения арифметических прогрессий $a+n\Z$. Она хаусдорфова ($\bigcap n\Z=0$; явно: при
$a\ne b$ и $n>|a-b|$ прогрессии $a+n\Z$ и $b+n\Z$ не пересекаются), но не дискретна ($\{0\}$ не открыто). Два целых
числа близки, если их разность делится на большое число.

\begin{center}
\begin{tikzpicture}[x=0.75cm]
\draw[->] (-6.6,0) -- (6.8,0);
\foreach \x in {-6,...,6} {\draw (\x,0.08)--(\x,-0.08); \node[below,font=\scriptsize] at (\x,-0.1) {$\x$};}
\foreach \x in {-6,-3,0,3,6} \fill[cold] (\x,0) circle (3pt);
\foreach \x in {-5,-2,1,4} \fill[cnew] (\x,0) circle (3pt);
\node[cold,font=\footnotesize] at (-4.5,0.5) {$0+3\Z$};
\node[cnew,font=\footnotesize] at (2.5,0.5) {$1+3\Z$};
\end{tikzpicture}
\end{center}
\figcap{Окрестности $0+3\Z$ и $1+3\Z$ точек $0$ и $1$ не пересекаются ($3>|0-1|$). С ростом $n$ окрестности $a+n\Z$
<<редеют>>, но остаются бесконечными~--- поэтому $\{0\}$ не открыто.}

\begin{center}
\begin{tikzpicture}[scale=0.9,every node/.style={font=\footnotesize}]
\draw[gray] (0,0) ellipse (4.4 and 2.3);
\node[gray] at (0,2.62) {$\Z=0+1!\,\Z=2+1!\,\Z$};
\draw[gray!60!black] (0,0) ellipse (3.55 and 1.78);
\node[gray!60!black] at (0,1.2) {$0+2!\,\Z=2+2!\,\Z$};
\foreach \c/\a/\col in {-1.6/0/cold,1.6/2/cnew} {
  \draw[\col,thick] (\c,-0.05) circle (0.9);
  \draw[\col] (\c,-0.05) circle (0.62);
  \draw[\col] (\c,-0.05) circle (0.38);
  \fill[\col] (\c,-0.05) circle (1.5pt) node[right=0.5pt,inner sep=0.8pt,font=\scriptsize]{$\a$};
  \node[\col] at (\c,-1.22) {$\a+3!\,\Z$};
}
\end{tikzpicture}
\end{center}
\figcap{Схема окрестностей точек $0$ и $2$. Окрестности $a+n!\,\Z$ ($n=1,2,3,\dots$) образуют базу в точке $a$ (любое
$m$ делит некоторое $n!$); внутренние круги~--- $a+4!\,\Z\supset a+5!\,\Z$. Вообще $b\in a+m\Z\iff m\mid(b-a)$, и два
класса по $m\Z$ либо совпадают, либо не пересекаются. Поэтому $0$ и $2$ лежат в общей окрестности $a+n!\,\Z$ при
$n\le2$, а при $n\ge3$ их окрестности не пересекаются: в большие окрестности любые две точки попадают вместе, в
достаточно малые~--- нет.}

\emph{Пополнение $\Zh$.} Последовательность $(a_i)$ целых чисел~--- \emph{последовательность Коши}, если для каждого
$n\ge1$ найдётся $N$ с $a_i\equiv a_j\pmod n$ при $i,j>N$; она \emph{тривиальна}, если для каждого $n$ найдётся $N$ с
$a_i\equiv0\pmod n$ при $i>N$. Последовательности Коши образуют коммутативное кольцо (покомпонентно), тривиальные~---
идеал в нём; фактор-кольцо обозначают $\Zh$. Отображение $m\mapsto(m,m,m,\dots)$ вкладывает $\Z$ в $\Zh$.

Удобнее другое описание. Последовательность Коши $(a_i)$ стабилизируется по модулю каждого $n$; обозначим инвариантный
вычет $c_n\in\Z/n$. Набор $(c_n)_{n\ge1}$ \emph{согласован}: $c_m\equiv c_n\pmod n$ при $n\mid m$. Последовательность тривиальна тогда
и только тогда, когда её набор нулевой (это просто переформулировка определения), и всякий согласованный набор
получается из некоторой последовательности Коши (возьмём
$a_k\equiv c_{k!}\pmod{k!}$: при $i,j\ge n$ число $n$ делит $i!$ и $j!$, и $a_i\equiv c_n\equiv a_j\pmod n$). Итак,
\[
\Zh=\{(c_n)_{n\ge1}\mid c_n\in\Z/n,\ c_m\equiv c_n\pmod n\text{ при }n\mid m\},
\]
с покомпонентными операциями~--- это \emph{обратный предел} $\varlim\Z/n\Z$ (\rf{7.20}--\rf{7.22}). Топология на $\Zh$~--- из
произведения дискретных $\Z/n$, база окрестностей нуля~--- $n\Zh=\{c\mid c_n=0\}$ (если $c_n=0$, то
$c_{nm}\in n\Z/nm\Z$ для всех $m$, т.\,е.\ $c=n\cdot c'$ для согласованного $c'$ с $c'_m=c_{nm}/n$); $\Z$ плотно в $\Zh$ (окрестность
$c+n\Zh$ содержит целое число~--- представитель $c_n$), а $\Zh$ компактно (замкнутое подмножество произведения
конечных множеств).

\emph{Изоморфизм.} Для $\alpha=(c_n)\in\Zh$ определим $\sigma^\alpha\in G$ условием
$\sigma^\alpha|_{\F_{p^n}}=\sigma_n^{c_n}$. Это корректно: при $n\mid m$ ограничение $\Gal(\F_{p^m}/\F_p)\to\Gal(\F_{p^n}/\F_p)$
переводит $\sigma_m$ в $\sigma_n$, т.\,е.\ соответствует редукции $\Z/m\to\Z/n$, и набор $(\sigma_n^{c_n})_n$ согласован; по \rf{7.8}
(шаг~4, или \rf{7.8a}(a)) он задаёт единственный элемент $G$. Отображение $\alpha\mapsto\sigma^\alpha\colon\Zh\to G$~---
изоморфизм групп (обратное: $\tau\mapsto(c_n)$, где $\tau|_{\F_{p^n}}=\sigma_n^{c_n}$) и гомеоморфизм (базисные
окрестности $n\Zh$ и $\Gal(\Om/\F_{p^n})$ соответствуют друг другу). При нём $1\mapsto\sigma$ и $\Z\mapsto\langle\sigma\rangle$.

\emph{Что это показывает.}
\begin{itemize}[leftmargin=*]
\item $\langle\sigma\rangle\cong\Z$~--- плотная, но не замкнутая подгруппа: $\Z\ne\Zh$. Например, элемент
$\alpha\in\Zh$ с $c_{2^k}=1$ и $c_{m}=0$ для нечётных $m$ (такой согласованный набор существует по китайской теореме об
остатках) не является целым числом; автоморфизм $\sigma^\alpha$ действует как Фробениус на всех $\F_{p^{2^k}}$ и
тривиально на всех $\F_{p^m}$ с нечётным $m$~--- никакая целая степень $\sigma$ так не действует. Более того, $\Zh$
несчётно (\rf{7.26b}).
\item Группа $\Zh$~--- главный пример \emph{проконечной группы} (\rf{7.24}). Для аналитиков она выглядит странно~--- её
компоненты связности одноточечны,~--- а для теории чисел естественна: по китайской теореме
$\Zh\cong\prod_{\ell\text{ простое}}\Z_\ell$, где $\Z_\ell$~--- кольцо $\ell$-адических целых чисел (\rf{7.33}).
\item Подрасширения $\Fpb$ и замкнутые подгруппы $\Zh$ описаны в \rf{7.34}.
\end{itemize}
\end{example}

\begin{example}[7.18: абсолютная группа Галуа поля $\Q$]
Пусть $\Qal$~--- алгебраическое замыкание $\Q$ в $\CC$. Группа $\Gal(\Qal/\Q)$~--- один из самых фундаментальных и
одновременно трудных объектов математики. Ожидается, что каждая конечная группа является её фактор-группой (обратная
задача теории Галуа); это известно, например, для $S_n$ и для всех спорадических простых групп, кроме, возможно, $M_{23}$.

Зато хорошо описывается группа $\Gal(F^{\rm ab}/F)$, где $F\subset\Qal$~--- конечное расширение $\Q$, а
$F^{\rm ab}$~--- объединение всех конечных абелевых расширений $F$ в $\Qal$; это предмет \emph{абелевой теории полей
классов}. Например, $\Gal(\Qab/\Q)\cong\Zh^\times$. Для $F=\Q$ ответ складывается из двух фактов:
\begin{itemize}[leftmargin=*]
\item $\Gal(\Q(\mu_\infty)/\Q)\cong\Zh^\times=\varlim(\Z/n)^\times$, где $\Q(\mu_\infty)$~--- поле, порождённое всеми
корнями из единицы (доказательство~--- в \rf{7.35}; оно использует лишь \rf{5.8} и \rf{5.10});
\item \emph{теорема Кронекера--Вебера} (без доказательства): каждое конечное абелево расширение $\Q$ содержится в
некотором $\Q(\mu_n)$, т.\,е.\ $\Qab=\Q(\mu_\infty)$.
\end{itemize}
Заметим, что $\Gal(\Qal/\Qab)$~--- замыкание коммутанта группы $\Gal(\Qal/\Q)$: фактор по замкнутой нормальной
подгруппе абелев тогда и только тогда, когда она содержит коммутант, и по \rf{7.13}(d) наибольшему абелеву подрасширению
отвечает наименьшая такая подгруппа. (В проконечной группе коммутант не обязан быть замкнут, поэтому замыкание здесь
нужно.)
\end{example}

\begin{aside}[7.19]
\emph{Простое соответствие Галуа}~--- это система из двух частично упорядоченных множеств $P$, $Q$ и обращающих порядок
отображений $f\colon P\to Q$, $g\colon Q\to P$ с $gf(p)\ge p$ и $fg(q)\ge q$. Тогда $fgf=f$: из $fg(fp)\ge fp$ и $gf(p)\ge
p$ следует $f(gfp)\le f(p)$. Аналогично $gfg=g$, и $f$, $g$ задают взаимно однозначное соответствие между $g(Q)$ и
$f(P)$.

Расширение Галуа $\Om/F$ даёт такую систему: $P$~--- подгруппы $\Gal(\Om/F)$, $Q$~--- подмножества $\Om$,
$f(H)=\Om^H$, $g(S)=G(S)$ (поточечный стабилизатор, теперь для любых $S$). Поэтому для взаимной однозначности в основной
теореме достаточно отождествить образ $g$ с замкнутыми подгруппами, а образ $f$~--- с промежуточными полями; это и
делает \rf{7.12}.
\end{aside}
